%%
%% 2019 07 04 Ph. G. Freimann
%%

\section{Potenzen}\index{Potenzen}
\sectuntertitel{Ein dreifaches Hoch auf die Basen}

\theorieGESO{57-59}{4.1, 4.2, 4.3 und 4.4}

%%%%%%%%%%%%%%%%%%%%%%%%%%%%%%%%%%%%%%%%%%%%%%%%%%%%%%%%%%%%%%%%%%%%%%%%%%%%%%%%%
\subsection*{Lernziele}

\begin{itemize}
\item Definition: Exponent, Basis, Potenz
\item Rechengesetze
\item Hoch 0
\item negative Exponenten
\item Potenzen von Bruchtermen
\item Namen der Zehnerpotenzen
\end{itemize}
\TALS{(\cite{frommenwiler17alg} S. 32 (Kap. 1.5))}
\newpage
%%\subsection{Definition}
%%$$a^n = \underbrace{a\cdot a \cdot a \cdot ... \cdot a}_{\textrm{n Faktoren}}$$


\subsection{Zehnerpotenzen}\index{Zehnerpotenzen}
\theorieTALS{32}{1.5.5}

\bbwCenterGraphic{8cm}{allg/alg/potenzen_wurzeln/img/one_in_a_mellon.jpg}

Neben den im Buch (\GESO{\cite{marthaler17}
  S. 62}\TALS{\cite{frommenwiler17alg} S. 40}) angegebenen SI-Einheiten (Kilo, Mega, ...) sind die
Namen der positiven Zehnerpotenzen im Englischen anders als im
Deutschen. Hier zur Vollstänigkeit:

\begin{tabular}{lclll}
Exponent   & SI-Kürzel & SI-Vorsätze & Deutsch & Englisch\\
\hline\\
$10^{2}$   & h         & Hekto       & Hundert   & hundred\\
$10^{3}$   & k         & Kilo        & Tausend   & thousand\\
$10^{6}$   & M         & Mega        & Million   & million\\
$10^{9}$   & G         & Giga        & Milliarde & billion\\
$10^{12}$  & T         & Tera        & Billion   & trillion\\
$10^{15}$  & P         & Peta        & Billiarde & quadrillion\\
$10^{18}$  & E         & Exa         & Trillion  & quintillion\\
\end{tabular}

\begin{tabular}{lclll}
Exponent   & SI-Kürzel & SI-Vorsätze & Deutsch\\
\hline\\
$10^{-1}$  & d         & Dezi        & Zehntel\\
$10^{-2}$  & c         & Centi       & Hundertstel\\
$10^{-3}$  & m         & Milli       & Tausendstel\\
$10^{-6}$  & $\mu$     & Mikro       & Millionstel\\
$10^{-9}$  & n         & Nano        & Milliardstel\\
\end{tabular}

\GESO{Mit dem Taschenrechner können große Zehnerpotenzen mit der \tiprobutton{EE}-Taste eingegben werden: 0.37 Milliarden: \tiprobutton{0}\tiprobutton{dot}\tiprobutton{3}\tiprobutton{7}\tiprobutton{EE}\tiprobutton{9}}%% END GESO
\newpage


\subsubsection{Ausklammern von Zehnerpotenzen}
Gegeben ist die folgende Summe. Leider etwas mühsam zum Lesen wegen der vielen Nullen. Klammern Sie Tausend (= $10^3$) aus:

$$5\,000 + 40\,000 + 6\,000\,000 + 70$$
\TNT{2.4}{$=10^3\cdot{} (5 + 40 + 6\,000 + 0.07 = 6054.07)$ \vspace{12mm}}


\paragraph{Ausklammern negativer Zehnerpotenzen:}
\,\\

Genauso, wie man positive Zehnerpotenzen ausklammern kann, kann man auch negative Zehnerpotenzen ausklammern. Dies ist insofern praktisch, um sich einen Überblick zu verschaffen, bei sehr kleinen positiven Größen:

Klammern Sie einen Millionstel (=$10^{-6}$) aus:


$$a\cdot{}10^{-6} + b\cdot{}10^{-2} + c\cdot{}10^{-5} + d\cdot{}10^{-1} = \LoesungsRaumLang{10^{-6} (a + b\cdot{}10^4 + 10c + d\cdot{}10^{5})}$$

\TALS{Zehnerpotenzen:}\TALSAadB{41ff}{101. a) b) c) d) 111. c) f) 112. b) d) e) f) k) m)}
\GESO{
  \subsection*{Aufgaben}
  \GESOAadB{65ff}{3. b) c), 4. c), 12., 15. und  16. a) und c)}
  \GESOAadB{68ff}{24. c), 28. c)}
  \GESOAadB{74ff}{52. a) c) d) 53. a) b) d) h) i) 56. a) b) 57. a) b) e) 60. 63. 64.}

}
\newpage


%% Philipp G Freimann Juli 2019 für die BBW
%% Phi BBW-Vorlage für Arbeitsblätter (LaTeX)
%% 2019 - 08 - 18

%% %% %% %%
\documentclass[twoside,12pt,a4paper]{article}%%

%%  In den Dokumenten können die folgenden Attribute überschrieben werden:

\newcommand{\metaHeaderLine}{HeaderLine mit $\backslash{}$ metaHeaderLine überschreiben}%%

%% \thema
\newcommand{\arbeitsblattTitel}{Pruefungsthema mit renewcommand arbeitsblattTitel überschreiben.}%%

\newcommand{\shortAuthor}{fp @ bbw.ch}

%%%%%%%%%%%%%%%%%%%%%%% P A C K A G E S %%%%%%%%%%%%%%%%%%%%%%%%%%%%%
\usepackage[paper=a4paper,margin=17mm]{geometry}%%


\usepackage{tgheros}%% Font TeX Gyre Heros für Titel (font code qhv)

%% damit die Punktezal schon geschrieben werden kann, obschon
%% Die Punkte erst während dem Dokument zusammengetragen werden:

%% Zentralisiert
\input{inputs/bbwUsePackages}
\input{inputs/bbwLayout}
\input{inputs/bbwMakros}
\input{inputs/matheMakros}
%%

%% LAYOUT FUER ARBEITSBLAETTER %%
\headheight12mm%%

%%%%%%%%%%%%%%%  H E A D E R   &   F O O T E R %%%%%%%%%%%%%%%%%%%%

%% Headers
\fancyhf[HL]{\makebox{\includegraphics[width=30mm]{logos/bbw.pdf}}}%%
\fancyhf[HC]{\metaHeaderLine{}}%%
\fancyhf[FR]{\tiny{\shortAuthor{} (\today{})}}%%

\newcommand{\arbeitsblattHeader}{%%
  \begin{center}%%
    {\Large \fontfamily{qhv}\selectfont \arbeitsblattTitel{}}%%
\end{center}}%%



%%%%%%%%%%%%%%%%%%%%%%%%%%%%%%%%%%%%%%%%%%%%%%%%%%%%%%%%%%%%%%%%%%

\usepackage{amssymb} %% für \blacktriangleright
\renewcommand{\metaHeaderLine}{Potenzgesetze}
\renewcommand{\arbeitsblattTitel}{(BMS)}

\begin{document}%%
\arbeitsblattHeader{}

\newcounter{aufgabennummer}
\setcounter{aufgabennummer}{1}

\newcommand\aufgabeML[2]{
Aufgabe \arabic{aufgabennummer}:\,\,
${#1}$ = $\LoesungsRaum{{#2}}$

\stepcounter{aufgabennummer}
}

{\huge{Vermischte Aufgaben zu Potenzgesetzen}}

\section{Ganze positive Exponenten}

\aufgabeML{(-10)^5}{-100\,000}
\aufgabeML{-10^7}{-10\,000\,000}
\aufgabeML{(-10)^8}{+100\,000\,000}
\aufgabeML{-(10^6)}{-1\,000\,000}
\aufgabeML{-10^4}{-10\,000}


\aufgabeML{-\left((-a)^3\right)^{10}}{-a^{30}}

\aufgabeML{\left(-(-x^2)\right)^7}{x^{14}}


\aufgabeML{\left(-x^3\right)^4}{x^{12}}
\platzFuerBerechnungen{14}
\noTRAINER{\newpage}


\section{Division}

\aufgabeML{c^9 : c^4}{c^5}

\aufgabeML{\frac{a^2}{a^3}}{\frac{1}{a}=a^{-1}}

\aufgabeML{(-a)^4:(-a^{10})}{-\frac{1}{a^6}}

\aufgabeML{a\cdot{} a^{2x+3} : a^{1-x}}{a^{3x+3}}

\aufgabeML{\left(\frac{x}{y}\right)^7:\left(\frac{-x}{y}\right)^3}{-\left(\frac{x}{y}\right)^4}
\platzFuerBerechnungen{18.8}
\noTRAINER{\newpage}

\section{negative Exponenten und die Null}

\aufgabeML{-\left(-(-a)^2\right)^0}{-1}

\aufgabeML{-\left( -a^3\cdot{} (-a)^2 \right)^{-4}}{\frac{-1}{a^{20}} = -a^{-20}}

\aufgabeML{(-b)^{-6} \cdot (-b^8)}{-b^2}

\aufgabeML{\left(-(a^{-1})^{-2}\right)^6}{a^{12}}

\aufgabeML{-\left((a^3)\cdot{}a^{-1}\right)^2}{-a^4}

Schreiben Sie ohne negative Exponenten:

\aufgabeML{2\cdot{}x^{-3}}{\frac{2}{x^3}}

\aufgabeML{a\cdot{}b^{-5}}{\frac{a}{b^5}}

\aufgabeML{\frac{1}{81}\cdot{}c^{-2}\cdot{}a^{-3}\cdot{}c^2\cdot{} a^3 \cdot 3^4}{1}
\platzFuerBerechnungen{16}
\noTRAINER{\newpage}

\section{negative Exponenten mit Brüchen}
\aufgabeML{\frac{6\cdot{} a^3 \cdot{} b^7 \cdot{} 3}{(ab)^4\cdot 9 \cdot a^{-1}}}{2b^3}

\aufgabeML{\left(\frac{a^2 b^{-3} c^3}{a\cdot b}\right)^{-3}\cdot \left(\frac{c^5}{ab}\right)^2}{\frac{b^{10}c}{a^5}}

\aufgabeML{\left(\frac{ab^{-2}c^3}{a^{-2}\cdot{}b}\right)^{-2} \cdot{} \left(\frac{a^3
\cdot{} c^{2}}{b^2} \right)^3}{a^3}

\aufgabeML{\left(\frac{a^3 \cdot
c^{-1}}{b^{-2}}\right)^{-4} \cdot \left( \frac{a^4\cdot{} b^3}{c^{-2}} \right)^3}{bc^{10}}

\aufgabeML{\frac{a^{-5}\cdot{} (-a)^3}{a^7} : \frac{a^{-10}}{a^4} }{-a^5}

\aufgabeML{\frac{ab\cdot{}b^{-2}\cdot{}b^4\cdot{}b^{-3}}{b^3\cdot{}
a\cdot{}b^{-4} \cdot{} b}}{1}

\aufgabeML{\left(\frac{3x^{-2}y^2}{4x^{-4}\cdot{}y^3}\right)^{-2} : \left(\frac{2x^{-1}}{3xy^{-2}}\right)^3}{\frac{6x^2}{y^4}}

\aufgabeML{4^k\cdot{} \left(\frac{1}{2}\right)^k \cdot{} \left(\frac{1}{3}\right)^{-k}}{6^k}

\aufgabeML{\frac{a^{-2}}{a^{-3}}}{a}

\aufgabeML{\left(\frac{a^4\cdot{}b^{-2}\cdot{}c}{a^2\cdot{}c^{-3}}\right)^{-2} \cdot \left(\frac{c^2\cdot{}b}{a^{-1}}\right)^4}{b^8\cdot{}c^4}

\aufgabeML{\left(\frac{3a^{-1}\cdot{} b^2}{2ac^{-1}}\right)^{-2} \cdot{} \frac{\left(3\cdot{}b^2\right)^2\cdot{}c^2}{4}}{a^4}

\aufgabeML{a^{-1}\cdot{} a^{-2} : a^{-3} \cdot{} a^4}{\frac{1}{a^4} = a^{-4}}

\platzFuerBerechnungen{12.4}
\noTRAINER{\newpage}



\section{Wurzeln}

\aufgabeML{\sqrt{a^2}}{a}

\aufgabeML{\left(\sqrt[2]{x}\right)^2}{x}

\aufgabeML{\sqrt[3]{b^3}}{b}

\aufgabeML{\sqrt[3]{z^9}}{z^3}

\aufgabeML{\sqrt[5]{b^{10}r^5}}{b^2r}

\aufgabeML{\sqrt[4]{r^8n^{12}}}{r^2n^3}

\aufgabeML{\sqrt[5]{m^4\cdot{m}}}{m}

\aufgabeML{\sqrt{m^2+m^2 + m^2}}{\sqrt{3}\cdot{}m}

\aufgabeML{\sqrt{\sqrt[2]{x^4}}}{x}

\aufgabeML{\frac{\sqrt[3]{a^7}}{\sqrt[3]{a}}}{a^2}

\aufgabeML{\sqrt[8]{r^{24}}}{r^3}

Und als Überleitung ins nächste Thema:

\aufgabeML{\sqrt[6]{v^3}}{v^{\frac12} = \sqrt{v}}

\platzFuerBerechnungen{12.8}
\noTRAINER{\newpage}




\section{Rationale Exponenten}

\aufgabeML{a^4\cdot{}a^{\frac{1}{4}}\cdot a^5\cdot a^{\frac{-1}{5}}}{a^\frac{181}{20}}

\aufgabeML{\sqrt[4]{a^3}}{a^{\frac{3}{4}}}

\aufgabeML{\sqrt[3]{r^{\frac{3}{4}}}}{\sqrt[4]{r} = r^{\frac{1}{4}}}

\aufgabeML{\sqrt{b^\frac12}}{\sqrt[4]b = b^\frac14}

\aufgabeML{\sqrt{a^\frac{1}{2}} \cdot \sqrt[4]{a\cdot{} \sqrt[5]{a^{10}}}}{a}

\aufgabeML{\sqrt{x^{10}} \cdot \sqrt[4]{x^3} \cdot x^{\frac{1}{4}}}{x^6}

\aufgabeML{\sqrt[3]{a} \cdot \sqrt{a^3 \cdot \sqrt[3]{a}}}{a^2}

\aufgabeML{\sqrt[3]{a^2} \cdot \sqrt[4]{a^3 \cdot \sqrt[3]{a}}}{a^{\frac{9}{6}}
= \sqrt[6]{a^9} = a^\frac32 = \sqrt{a^3}}

\aufgabeML{\sqrt[3]{b\cdot \sqrt{b^3\cdot \sqrt[3]b}}}{\sqrt[9]{b^8}}

\aufgabeML{\sqrt[4]{b\cdot{} \sqrt[3]{b^2\cdot \sqrt{b}}}}{b^\frac{11}{24}
= \sqrt[24]{b^{11}}}

\aufgabeML{3a\cdot \sqrt[3]{9a^2}}{\sqrt[3]{3^5 a^5} = \left(3a\right)^\frac53}

\aufgabeML{\sqrt{x} \cdot \sqrt[3]{x^4} \cdot \sqrt[6]{x^3}}{x^\frac73
= \sqrt[3]{x^7}}

\aufgabeML{\sqrt{a^3}\cdot \sqrt[3]{a^2}}{a^\frac{13}{6} = \sqrt[6]{a^{13}}}

\aufgabeML{\sqrt{a^2 \sqrt{a}}}{a^\frac54 = \sqrt[4]{a^5}}

\aufgabeML{\sqrt{a^{-2}}}{\frac{1}{a} = a^{-1}}

\aufgabeML{\sqrt[3]{a^2 \cdot b \cdot \sqrt{a\cdot b^{-1}}}}{a^\frac56
\cdot b^\frac16 = \sqrt[6]{a^5\cdot b}}

\aufgabeML{\sqrt{x^3} \cdot \left(\sqrt{x}\right)^{-3}}{1}

\aufgabeML{\frac{\sqrt{a^3}}{\sqrt{a^{-1}}}}{a^2}

\aufgabeML{\left(\frac{1}{a}\right)^{-\frac14}}{\sqrt[4]a = a^\frac14}

\aufgabeML{\frac{a^4}{\sqrt{a}}}{\sqrt{a^7} = (\sqrt{a})^7 = a^\frac72}

\aufgabeML{\frac{\sqrt[3]{a^2}}{\sqrt{a}}}{a^{\frac16} = \sqrt[6]{a}}

\aufgabeML{\frac{-\sqrt{a^3}}{-\left(\sqrt{a}\right)^3}}{1}

\aufgabeML{\frac{\sqrt[3]{a^{13}}}{a^4}}{a^{\frac13} = \sqrt[3]{a}}
\noTRAINER{\newpage}

\platzFuerBerechnungen{23.2}

\end{document}


\subsubsection{Negative Exponenten}
\sectuntertitel{Nicht für alle ist die Potenzrechnung \textit{positiv}.}
Wir kennen bereits das Rechengesetz für positive Exponenten:

$$a^5\cdot{}a^2 = a^{5+2} = a^7$$

Sinnvoll wäre die folgende Erweiterung auf negative
Exponenten:\\\TRAINER{Damit die Rechengesetze weiterhin gelten.}
$$a^5\cdot{}a^{-2} = a^{5+(-2)} = a^3$$

Dividieren wir obige Gleichung beidseitig durch $a^5$, so erhalten wir
folgende sinnvolle Definition:

\TNT{3.2}{
Wir wollen erreichen, dass gilt: $a^5\cdot{}a^{-2} \stackrel{!}{=}a^3$

Wir dividieren durch $a^5$ beidseitig: $a^{-2} = \frac{a^3}{a^5} = \frac{1}{a^2}$

Mit dieser Herleitung von $a^{-2}= \frac1{a^2}$ gilt: $a^5\cdot{}a^{-2} = a^5\cdot{}\frac1{a^2} = a^3$
}%% END TNT

\begin{definition}{}{}
$$a^{-n} := \frac{1}{a^n}$$
\end{definition}

\begin{bemerkung}{}{}
$$\frac{1}{a^n}= 1 : \underbrace{a : a : a : ... : a}_{n \textrm{\ Divisoren}}$$
\end{bemerkung}

\begin{gesetz}{}{}
$a^{-n} = \left(\frac1a\right)^n$
\end{gesetz}
Begründung:

\TNT{4.4}{$a^{-n} = \frac1{a^n}
= \frac1{\underbrace{a\cdot{}a\cdot{}a\cdot{}...\cdot{}a}_{n \textrm{\ Faktoren}}}
= \underbrace{\frac1a\cdot{}\frac1a\cdot{}\frac1a\cdot{}...\cdot{}\frac1a}_{n \textrm{\ Faktoren}}
= \left(\frac1a\right)^n$
}%% END TNT
\newpage
Ganz analog gilt:

\begin{gesetz}{}{}
$\frac{1}{a^{-n}} =a^n$
\end{gesetz}
Begründung:
\TNT{8}{


Beispiel:

$$\frac1{10^{-3}} = \frac1{0.001} = 1000 = 10^3$$

\TALS{TALS:}
Beweis: Definition hinschreiben und auf beiden Seiten den Kehrwert bilden:

$$a^{-n} = \frac1{a^n}$$

$$\frac1{a^{-n}} = a^n$$

\vspace{2cm}
}%% END TNT



\begin{gesetz}{}{}
$\left(\frac{1}{a}\right)^{-n}=a^n$
\end{gesetz}
Begründung:
\TNT{8}{
Zahlenbeispiel (erster Schritt nach Definition):
$$\left(\frac1{10}\right)^{-3}  = \frac1{\left(\frac1{10}\right)^3}  = \frac1{0.001} = 1000 = 10^3$$

\TALS{TALS:}

Beweis: Nach Definition gilt für alle $n$:

$$x^{-n} = \frac1{x^n}$$

Somit gilt es auch, wenn wir anstelle von $x$ den Term $\frac1a$ einsetzen:

$$\left(\frac1a\right)^{-n} = \frac1{\frac1{a^n}} = a^n$$


\vspace{2cm}

}%% END TNT




\newpage
Rechenbeispiel:

Wurm «Wurli» schaft 3 cm pro Sekunde (= $3 \cdot{} 10^{-2} $ m pro Sekunde). Wie lange braucht «Wurli» für 12 m?
\TNT{2.4}{
$$t = \frac{s}v = \frac{12[ \textrm{m}]}{3 \frac{[\textrm{cm}]}{[\textrm{s}]}} =   \frac{12[ \textrm{m}]}{3\cdot{}10^{-2}\frac{[\textrm{m}]}{[\textrm{s}]}} = 4\cdot{} 10^2 [\textrm{s}] = 400 [\textrm{s}]\approx 6-7 Min.$$
}%% END TNT




Und ebenso für beliebige Brüche:

\begin{gesetz}{}{}
$\left(\frac{a}{b}\right)^{-n} = \left(\frac{b}{a}\right)^{+n}$
\end{gesetz}

    
    \TALS{ Begründung

      \TNT{2.4}{$\left( \frac{b}{a} \right)^n  =
       \left(b \cdot{} \frac{1}{a} \right)^n =
       b^n \cdot \left(\frac{1}{a}\right)^n =
       \left(\frac{1}{b}\right)^{-n} \cdot{} a^{-n} =
       \left(\frac{1}{b}\cdot{}a\right)^{-n} =
       \left(\frac{a}{b}\right)^{-n} 
       $}} %% END TNT END TALS

    \GESO{ Begründung

      \TNT{2.4}{$\left( \frac{5}{2} \right)^3  =
       \left(5 \cdot{} \frac{1}{2} \right)^3 =
       5^3 \cdot \left(\frac{1}{2}\right)^3 =
       \left(\frac{1}{5}\right)^{-3} \cdot{} 2^{-3} =
       \left(\frac{1}{5}\cdot{}2\right)^{-3} =
       \left(\frac{2}{5}\right)^{-3} 
       $}} %% END TNT END GESO

%%$$\left(\frac{1}{a}\right)^{-n} = \frac{1}{\left(\frac{1}{a}\right)^n} = a^n$$
\newpage



\subsubsection{Null}

Ebenso müsste $a^5 \cdot a^0 = a^{5+0} = a^5$ gelten. Dies ist aber
nur möglich, wenn wir $a^0 := 1$ definieren ($a\ne 0$).

\begin{definition}{Exponent Null}{} Für alle Basen
$a \in \mathbb{R}\backslash\{0\}$ definieren wir:
\begin{center}
\fbox{$a^0 := 1$}
\end{center}
\end{definition}

Zweite Begründung wenn die Rechengesetze gelten sollten: $a^0 = a^{1-1} = \frac{a^1}{a^1} = 1$

%\textbf{Rechengesetze zusammengefasst:}

%\begin{itemize}
%\item  $\frac{a^m}{a^n} = a^{m-n}$ (Dies gilt auch wenn $n > m$.)
 
%\item $a^{-n} := a^{0-n}=\frac{a^0}{a^n} = \frac{1}{a^n} = \left(\frac{1}{a}\right)^n$

%\item
%$\left(\frac{1}{a}\right)^{-n} = \frac{1}{\left(\frac{1}{a}\right)^n} = a^n$


%\item $\left(\frac{a}{b}\right)^{-n} = \left(\frac{b}{a}\right)^{+n}$ gilt daher auch. 
%\end{itemize}

\begin{rezept*}{«Kielholen»}{}{}
Exponenten vertauschen ihr Vorzeichen beim Übertreten des Bruchstrichs:
$$\frac{a^{-3}b^2}{c^5d^{-6}} = \frac{b^2d^6}{a^3c^5}$$
\end{rezept*}



\subsection{Aufgaben}

\TALS{Potenzen:}\TALSAadB{32ff}{Von Hand: 79. c), 82. a), 83. b),
86. b) c),
91. a), 92. j), 94. k), 96. b) f) h) und 105. i)\\
Prüfen Sie die folgenden Aufgaben auch mit dem TR (Training):\\
103. a) b) c), und 106. h}

\GESOAadB{67ff}{15., 18. c), 19. b), 20. h), 26. b), 31. b),
  38. c) e), 41. e), 43. c), 44. d) e) f) h) i), 48. a) b), 49. a) c)}

\aufgabenfarbe{Blatt zu Potenzgesetzen im \texttt{olat.bbw.ch}; dort Kap. 1.3 und 1.4}

\GESO{Optional: \GESOAadB{72ff}{51. (Koch)}}

\newpage

